% example.tex
% 编译方式：xelatex example.tex

\documentclass[aspectratio=169,11pt]{beamer}

\usepackage{fontspec}
\setsansfont{Cochin}
\usepackage{xeCJK}
\setCJKmainfont{Songti SC}

\usetheme{Boadilla}
\usepackage{../style/titlefoot}  % 必须在主题后面才生效

% ===== 元信息 =====
\title[短标题：Beamer 页脚自动切换]{Beamer 页脚自动切换示例}
\subtitle{正文 / 标题页 / 目录页 三种页脚}
\author[张三]{张三\inst{1} \and 李四\inst{2}}
\institute[XX 大学]{
  \inst{1} XX 大学\quad XX 学院 \\
  \inst{2} XX 大学\quad XX 学院
}
\date[2026-10-10]{2026 年 10 月 10 日}

\begin{document}

% =========================================================
% 1. 标题页：应只显示 short title
% =========================================================
\begin{frame}
  \titlepage
\end{frame}

% =========================================================
% 2. 目录页：应只显示 short title
% =========================================================
\begin{frame}
  \frametitle{目录}
  \tableofcontents
\end{frame}

% =========================================================
% 3. 第一节：正文页脚应显示 author / section / 页码
% =========================================================
\section{第一节：基本用法}

\begin{frame}
  \frametitle{普通正文页}
  这一页的页脚应显示：
  \begin{itemize}
    \item 左侧：\texttt{\textbackslash insertshortauthor}，即“张三”
    \item 中间：\texttt{\textbackslash insertsection}，即“第一节：基本用法”
    \item 右侧：\texttt{\textbackslash insertframenumber/\textbackslash inserttotalframenumber}
  \end{itemize}
\end{frame}

\begin{frame}
  \frametitle{带列表的正文页}
  \begin{enumerate}
    \item 第一项
    \item 第二项
    \item 第三项
  \end{enumerate}
\end{frame}

\begin{frame}
  \frametitle{带公式的正文页}
  \begin{equation}
    E = mc^2
  \end{equation}
  \[
    \int_{-\infty}^{+\infty} e^{-x^2}\,\mathrm{d}x = \sqrt{\pi}
  \]
\end{frame}

% =========================================================
% 4. 第二节：验证 section 自动切换
% =========================================================
\section{第二节：页脚切换验证}

\begin{frame}
  \frametitle{进入第二节后的页脚}
  此时中间栏应自动变为“第二节：页脚切换验证”。
\end{frame}

\begin{frame}
  \frametitle{再一页}
  继续验证页码递增。
\end{frame}

% =========================================================
% 5. 第三节：分栏、块环境等复杂内容
% =========================================================
\section{第三节：复杂内容}

\begin{frame}
  \frametitle{分栏布局}
  \begin{columns}[T]
    \begin{column}{0.48\textwidth}
      \begin{block}{左栏}
        左侧内容。页脚仍然正常显示。
      \end{block}
    \end{column}
    \begin{column}{0.48\textwidth}
      \begin{block}{右栏}
        右侧内容。
      \end{block}
    \end{column}
  \end{columns}
\end{frame}

\begin{frame}[fragile]
  \frametitle{代码页}
\begin{verbatim}
\setbeamertemplate{footline}{
  \insertshortauthor \hfill \insertsection \hfill
  \insertframenumber/\inserttotalframenumber
}
\end{verbatim}
\end{frame}

\begin{frame}
  \frametitle{警告与示例块}
  \begin{alertblock}{注意}
    标题页和目录页的页脚只显示短标题。
  \end{alertblock}
  \begin{exampleblock}{示例}
    正文页页脚分为三栏：作者、节名、页码。
  \end{exampleblock}
\end{frame}

% =========================================================
% 7. 结束页
% =========================================================
\begin{frame}
  \thispagestyle{navigation@titlepage} % 使用首页页脚样式
  \centering
  \Huge 谢谢！
\end{frame}

\end{document}